% !TEX root = ../matan.tex

\section{Условно сходящиеся ряды, сравнение с абсолютно сходящимися, ряд Лейбница}

\subsection*{Абсолютная и условная сходимость}

\begin{Statement}{Абсолютная сходимость влечёт сходимость}{}
	Если $\sum_{k=1}^{\infty} |a_k|$ сходится, то и $\sum_{k=1}^{\infty} a_k$
	сходится, причём
	$\bigl|\sum_{k=1}^{\infty} a_k\bigr|\le\sum_{k=1}^{\infty} |a_k|$.
\end{Statement}

\begin{proof}
	Применим критерий Коши. Из сходимости $\sum|a_k|$: для любого $\eps>0$
	найдётся $N$ такое, что при $n\ge N$ и любом $m$ выполнено
	$\sum_{k=n+1}^{n+m}|a_k|<\eps$. Тогда по неравенству треугольника
	$\bigl|\sum_{k=n+1}^{n+m} a_k\bigr|\le\sum_{k=n+1}^{n+m}|a_k|<\eps$, то есть
	для $\sum a_k$ выполнен критерий Коши и ряд сходится. Оценка суммы получается
	предельным переходом из $|S_n|\le\sum_{k=1}^n|a_k|$.
\end{proof}

\begin{Remark}{Сравнение двух типов сходимости}{}
	Обратное неверно: условная сходимость абсолютной не влечёт. Принципиальное
	различие проявляется при \emph{перестановке} членов: сумма абсолютно
	сходящегося ряда не зависит от порядка слагаемых, тогда как у условно
	сходящегося ряда перестановкой членов можно получить любую наперёд заданную
	сумму (теорема Римана). Таким образом, абсолютная сходимость --- это
	«безусловная» сходимость, не чувствительная к порядку, а условная
	сходимость существенно от порядка зависит.
\end{Remark}

\subsection*{Ряд Лейбница}

\begin{Definition}{Ряд Лейбница (знакочередующийся ряд)}{}
	Ряд $\sum_{k=1}^{\infty} a_k$ называется \textbf{рядом Лейбница}, если
	\begin{enumerate}
		\item знаки соседних членов чередуются:
		$\operatorname{sign} a_{k+1}=-\operatorname{sign} a_k$, $a_k\ne0$;
		\item модули монотонно убывают: $|a_{k+1}|\le|a_k|$;
		\item $|a_k|\to0$ при $k\to\infty$.
	\end{enumerate}
	Такой ряд можно записать в виде $\sum_{k=1}^{\infty} (-1)^{k-1} c_k$, где
	$c_k:=|a_k|$, $c_k\searrow0$.
\end{Definition}

\begin{Theorem}{Признак Лейбница}{}
	Всякий ряд Лейбница сходится.
\end{Theorem}

\begin{proof}
	Вынося общий знак, можно считать $a_k=(-1)^{k-1} c_k$ с $c_k\searrow0$,
	$c_k\ge0$. Рассмотрим произвольный хвост, начинающийся с индекса $n+1$:
	для $m\ge1$
	\[
	T_m:=\sum_{k=n+1}^{n+m} a_k
	=(-1)^{n}\bigl(c_{n+1}-c_{n+2}+c_{n+3}-\cdots\pm c_{n+m}\bigr).
	\]
	Обозначим $U_m:=c_{n+1}-c_{n+2}+\cdots\pm c_{n+m}$ (так что $|T_m|=|U_m|$).
	Сгруппируем слагаемые двумя способами. Группируя с начала парами,
	\[
	U_m=(c_{n+1}-c_{n+2})+(c_{n+3}-c_{n+4})+\cdots\ge0,
	\]
	так как каждая скобка неотрицательна ($c_k$ убывает); незакрытый последний
	член (если $m$ нечётно) также неотрицателен. С другой стороны,
	\[
	U_m=c_{n+1}-(c_{n+2}-c_{n+3})-(c_{n+4}-c_{n+5})-\cdots\le c_{n+1},
	\]
	так как все вычитаемые скобки неотрицательны. Итак, $0\le U_m\le c_{n+1}$, то
	есть
	\[
	|T_m|=\Bigl|\sum_{k=n+1}^{n+m} a_k\Bigr|\le c_{n+1}=|a_{n+1}|
	\qquad\forall m\ge1.
	\tag{$\ast$}
	\]
	Так как $c_{n+1}\to0$ при $n\to\infty$, для любого $\eps>0$ найдётся $n_0$
	такое, что при $n\ge n_0$ имеем $|T_m|\le c_{n+1}<\eps$ сразу для всех $m$.
	По критерию Коши ряд сходится.
\end{proof}

\begin{Example}{Знакочередующийся гармонический ряд}{}
	Ряд $\displaystyle\sum_{k=1}^{\infty}\frac{(-1)^{k-1}}{k}
	=1-\frac12+\frac13-\frac14+\cdots$ сходится по признаку Лейбница, так как
	$c_k=\frac1k$ монотонно убывает к нулю. Однако он \textbf{не} сходится
	абсолютно: ряд модулей $\sum\frac1k$ --- гармонический --- расходится.
	Следовательно, это пример \textbf{условно} сходящегося ряда. (Его сумма
	равна $\ln2$.)
\end{Example}
